\subsection{Hello Again}
\setauthor{Zoe Emily Öllinger}

Hello Again ist ein österreichisches Softwareunternehmen, das sich auf digitale Kundenbindungslösungen spezialisiert hat.
Das Kernprodukt ist eine Kundenbindungsapp, mit der Unternehmen ihren Kunden digitale Stempelkarten, Punkte und Prämien sowie Gewinnspiele anbieten können.
Die App wird von Hello Again entwickelt und betrieben, aber im Namen der jeweiligen Kundenfirmen veröffentlicht.


Das Unternehmen betreut viele Kunden gleichzeitig.
Die Apps sind sowohl im Google Play Store als auch im Apple App Store verfügbar.
Um den Überblick über den Zustand aller Apps und Kundenaccounts zu behalten, wurde das Health Monitoring Dashboard entwickelt.


\subsection{Projektteam und Zusammenarbeit}
\setauthor{Zoe Emily Öllinger}

Wir waren als Praktikanten zu Hello Again gekommen und wurden den Devs vorgestellt.
Zu Beginn des Projekts fand eine Einschulung statt, in der uns die wichtigsten Richtlinien und Tools erklärt wurden.
Die Betreuung erfolgte durch einen fixen Ansprechpartner im Unternehmen, an den wir uns bei Fragen oder Problemen wenden konnten.


Wir arbeiteten in einem eigenen Git-Branch namens health-monitoring, der vom offiziellen Main-Branch getrennt war.
Innerhalb dieses Branches haben wir uns zwei weitere Branches angelegt einen für das Frontend und einen für das Backend.
Am Ende des Projekts wurde unsere Arbeit von den Entwicklern geprüft und in das offizielle Projekt übernommen.


\subsection{Markt und Branche}
\setauthor{Zoe Emily Öllinger}

Der Markt für Kundenbindungslösungen wächst seit einigen Jahren kontinuierlich.
Immer mehr Unternehmen ersetzen ihre physischen Stempelkarten durch digitale Apps, wie zum Beispiel Müller, weil diese einfacher zu verwalten sind und gleichzeitig Daten über das Kaufverhalten der Kunden liefern.





